\part{Pre-training}

\chapter{Pre-training Objectives and Strategies}
\minitoc

\begin{chapteroverview}
  \begin{itemize}
    \item Master the next-token prediction (NTP) objective and its 2026 variants.
    \item Review Reinforcement Pre-Training (RPT) for enhanced reasoning.
    \item Analyze curriculum learning and instruction-augmented pre-training.
    \item Understand the multi-stage compute-optimal context extension pipeline.
  \end{itemize}
\end{chapteroverview}

\section{Next-Token Prediction}
Self-supervised causal language modeling (see also \hyperref[form:ntp]{\S\ref*{form:ntp}} for the full NTP objective and \hyperref[form:ppl]{\S\ref*{form:ppl}} for perplexity (Appendix~\ref{app:objectives}); code: \hyperref[lst:ntp_loss]{Listing~\ref*{lst:ntp_loss}}):
\[
  \mathcal{L} = -\sum_{t=1}^{T} \log p_\theta(x_t \mid x_1, \ldots, x_{t-1})
\]
No labeled data needed. Captures syntax, semantics, factual knowledge, and reasoning patterns. Despite its simplicity, this objective is sufficient to produce emergent capabilities at scale.

\section{Reinforcement Pre-Training (RPT)}
Microsoft (2025): next-token prediction reframed as a sequential decision-making process. Each token prediction is a policy action; the model receives richer gradient signals than maximum likelihood. Fully self-supervised. Improves long-horizon coherence and reasoning traces.

\section{Curriculum Learning}
Data organized by difficulty (simple $\rightarrow$ complex) accelerates convergence and boosts final performance. Difficulty metrics: perplexity under a smaller reference model, task complexity (number of reasoning steps required), domain specificity.

\section{Instruction-Response Augmented Pre-training}
Synthetic instruction-response pairs in the pre-training corpus bridge the gap to SFT. Typically 1--3\% of the pre-training mix. Reduces the SFT data required by 5--10$\times$ for comparable instruction-following capability.

\section{Long-Context Pre-training}

Short-context training followed by gradual context extension:
\begin{enumerate}
  \item Train at 4K context for 90\% of total compute.
  \item Gradually increase to 32K--128K on a filtered long-document subset.
  \item Apply YaRN or LongRoPE for positional interpolation.
\end{enumerate}
Training from scratch at long context is compute-inefficient; the model wastes FLOPs on mostly short sequences in the dataset.
